\documentclass[11pt]{article}
\usepackage[margin=1in]{geometry}
\usepackage{amsmath,amssymb}
\usepackage{booktabs}
\usepackage{xcolor}
\usepackage[hidelinks]{hyperref}

\newcommand{\code}[1]{[\![#1]\!]}
\newcommand{\dleq}{d\!\le}
\newcommand{\tablecmd}[1]{\texttt{#1}}

% Let TeX break before long unbreakable \texttt identifiers instead of
% overflowing the margin (e.g. scipy.optimize.milp, mip_dual_bound).
\setlength{\emergencystretch}{3em}

\title{Independent cross-check of quantum LDPC code distances against the
\texttt{codeDistance} package, and a route to distance bounds for larger codes}
\author{Juan Cruz-Benito}
\date{\today}

\begin{document}
\maketitle

\begin{abstract}
We cross-checked the code-distance values recorded by our discovery
pipeline~\cite{qcode} (BP-OSD plus a SciPy/HiGHS MILP) against the independent
\texttt{codeDistance} package of
Webster, Jacob, and Higgott (arXiv:2603.22532). Across 198 distinct codes (43 CSS,
155 non-CSS perturbed-bicycle), independent methods matched or corroborated our
recorded values with \emph{zero contradictions}, including the codes where BP-OSD
overestimated the distance by up to $12\times$. We then show, on Webster's own
evolutionary search \texttt{QDistEvol}, that a cheap heuristic matches the recorded
MILP value on every one of the worst overestimate cases and even tightens one of
our incumbents. Finally we analyse what these methods can and cannot do for
\emph{larger} codes, where the MILP times out, and give a concrete
bound-bracketing workflow that combines a heuristic upper bound with the certified
lower bound HiGHS can report but our current pipeline does not persist.
\end{abstract}

\section{Setup and protocol}

The codes under test are the bivariate-bicycle (BB) and perturbed-bivariate-bicycle
(PBB) quantum LDPC codes discovered in~\cite{qcode}, whose distances were recorded
there via BP-OSD and a SciPy/HiGHS MILP (catalogs \texttt{ilp\_catalog.json} for
CSS and the campaign-7 record for non-CSS).

The \texttt{codeDistance} package exposes \texttt{CSScodeDistance} (CSS) and
\texttt{codeDistance} (non-CSS, symplectic two-block input). We drive four
methods: \texttt{BZDistMW} (Brouwer--Zimmermann enumeration, exact on
completion), \texttt{MIPDist} (OR-Tools/SCIP), \texttt{QDistEvol} (evolutionary
heuristic), and \texttt{decoderDist} (BP-OSD via the \texttt{ldpc} package).

The protocol rests on one principle: \textbf{trust the returned witness, not the
reported distance}. For every run we recompute the weight of the returned operator
(Hamming for CSS, symplectic for non-CSS) and verify that it commutes with all
checks and anticommutes with at least one logical. A \emph{verified} witness of
weight $w$ is then a hard certificate that $\dleq w$, independent of the method and
of solver bookkeeping. Two consequences shape the reading of every result below:

\begin{itemize}
  \item \texttt{MIPDist} is always recorded as an \emph{incumbent}: the public
  package result does not expose OR-Tools' optimality status, so we never promote
  a finished solve to a proof. \texttt{BZDistMW} is \texttt{proven\_exact} only
  when it runs to completion.
  \item The highest-impact signal is a verified witness \emph{lighter} than a
  distance we claim exact (the \texttt{CRITICAL\_LOWER} verdict). None was
  observed.
\end{itemize}

\section{Results}

\paragraph{Overall.} 198 distinct codes, 271 method-runs (254 \texttt{MIPDist}
incumbents, 17 \texttt{BZDistMW} proofs), every returned witness verified.
Table~\ref{tab:verdicts} gives the per-code verdict.

\begin{table}[h]
\centering
\begin{tabular}{lcl}
\toprule
verdict & codes & meaning \\
\midrule
\textbf{agree} & \textbf{196} & min verified witness weight $=$ our recorded value \\
\texttt{webster\_higher} & 2 & heavier valid logical (looser bound), not a contradiction \\
\textbf{CRITICAL} & \textbf{0} & no verified witness below a claimed-exact distance \\
\bottomrule
\end{tabular}
\caption{Per-code verdicts over the 198-code cross-check. The two
\texttt{webster\_higher} rows are $n{=}108$, $d{=}10$ codes whose MIP returned
weight 12 even at a 600\,s budget ($12 \ge 10$, consistent).}
\label{tab:verdicts}
\end{table}

\paragraph{CSS (Comparison A).} All 43 CSS codes agree: the two Bravyi baselines
($\code{72,12,6}$, $\code{144,12,12}$), 32 catalog rows at $n{=}144$, and the 9
``BP-OSD overestimated'' rows at $n{=}288/360$. We stress that
\texttt{ilp\_catalog.json} labels are not solver certificates; the non-Bravyi rows
are \emph{corroborated} by independent witnesses, not upgraded to proven-exact.

The headline cross-check is the overestimate subset (Table~\ref{tab:overest}):
these are the codes where our pipeline most aggressively corrected BP-OSD. An
independent MIP (OR-Tools/SCIP, distinct from our SciPy/HiGHS formulation)
returned a verified logical at our recorded value in both $X$ and $Z$ on all
nine. Because a verified weight-$w$ witness certifies $\dleq w$, this
\emph{proves} the upper bounds that refute the BP-OSD estimates---e.g.
$\code{360,40}$ with a weight-2 witness against a BP-OSD estimate of 24, a
$12\times$ overestimate. Equality with our recorded value is corroboration, not an
independent optimality proof (all runs are incumbents, and these catalog rows are
incumbents too); no lighter witness was ever found.

\begin{table}[h]
\centering
\begin{tabular}{lccc}
\toprule
code & BP-OSD $d$ & MILP / verified witness & BP-OSD\,/\,witness \\
\midrule
$\code{360,40}$ & 24 & \textbf{2} & $\mathbf{12.0\times}$ \\
$\code{360,40}$ & 22 & \textbf{2} & $11.0\times$ \\
$\code{360,40}$ & 20 & \textbf{2} & $10.0\times$ \\
$\code{288,32}$ & 22 & \textbf{4} & $5.5\times$ \\
$\code{360,32}$ & 20 & \textbf{4} & $5.0\times$ \\
$\code{360,32}$ & 16 & \textbf{4} & $4.0\times$ \\
$\code{288,32}$ & 20 & \textbf{6} & $3.3\times$ \\
$\code{288,32}$ & 12 & \textbf{4} & $3.0\times$ \\
$\code{288,24}$ & 18 & \textbf{12} & $1.5\times$ \\
\bottomrule
\end{tabular}
\caption{The nine ``BP-OSD overestimated'' CSS codes. The verified witness weight
matches our recorded MILP value on all nine, certifying the corrected (much lower) upper
bound against BP-OSD overestimates of up to $12\times$.}
\label{tab:overest}
\end{table}

\paragraph{Non-CSS (Comparison C).} Of 155 exact non-CSS codes ($n\le180$,
$d>6$, plus the $n{=}36$ baselines), 153 agree and 2 are looser; none
contradicts. The non-CSS symplectic linearization our MILP uses therefore held on
every code tested: every returned witness was a valid logical and none was lighter
than our recorded distance. A targeted 600\,s rerun of an initial 13 looser
incumbents closed 11 of them to agreement with our exact record.

\paragraph{Heuristics (Comparison B): does \texttt{QDistEvol} close the BP-OSD
gap?} On the 17 worst overestimate codes (the 9 CSS rows plus the 8 largest
non-CSS $d_{\text{BP-OSD}}-d_{\text{MILP}}$ gaps) we ran \texttt{QDistEvol} and
\texttt{decoderDist} and compared each verified witness weight to our recorded
MILP value (Table~\ref{tab:tier2}). \texttt{QDistEvol} reached the recorded MILP value on
\emph{all 17}, and on one non-CSS code ($\code{360,12}$, BLISS
\texttt{c447dcb2d101746f}) it \emph{improved} our incumbent, finding a weight-20
logical where the recorded value was 24. It is cheap ($\sim$46--133\,s/code).

Two nuances sharpen the BP-OSD story. On CSS, \texttt{decoderDist} (a well-run
BP-OSD) also matches the recorded \texttt{ilp\_d} value on all nine---so the
overestimates recorded in our catalog reflect \emph{our pipeline's} BP-OSD
configuration, not an intrinsic BP-OSD limitation. On non-CSS, \texttt{decoderDist}
timed out with no verified logical on all eight; \texttt{QDistEvol} was the only
Tier-2 heuristic that returned valid witnesses there. This is exactly the regime
where a heuristic alternative to BP-OSD matters.

\begin{table}[h]
\centering
\begin{tabular}{lcc}
\toprule
set & \texttt{decoderDist} reaches recorded MILP & \texttt{QDistEvol} reaches recorded MILP \\
\midrule
CSS (9) & 9/9 & 9/9 \\
non-CSS (8) & 0/8 (timed out, no logical) & 8/8 \\
\textbf{total} & \textbf{9/17} & \textbf{17/17} \\
\bottomrule
\end{tabular}
\caption{Gap-closing on the worst BP-OSD-overestimate codes.}
\label{tab:tier2}
\end{table}

\section{Distance bounds for larger codes}

The pressing practical problem is that for most $n\ge108$ codes the MILP times out
before proving optimality. We separate two things a timeout costs us, because the
available remedies differ:
\[
\text{upper bound: } \dleq U \quad\text{(a low-weight logical)},\qquad
\text{lower bound: } d\ge L \quad\text{(a proof nothing is lighter).}
\]

\paragraph{Upper bounds are sound at any size; the limit is tightness.} Since we
verify every witness, ``method returned a weight-$w$ logical $\Rightarrow \dleq w$''
holds for arbitrarily large codes---the verification is $\mathrm{GF}(2)$ linear
algebra. What degrades with size is whether a method \emph{finds} a tight witness.
Here \texttt{QDistEvol} genuinely helps: it is cheap, independent, and on our
hardest non-CSS codes it matched the MILP and even beat an incumbent
(Table~\ref{tab:tier2}). \texttt{MIPDist} does \emph{not} break the wall---it is
the same NP-hard problem in a different solver and returned only incumbents on the
hard codes---so its value at scale is a tight \emph{witness}, not a certificate.

\paragraph{The lower bound is not being persisted.} Certifying $d=w$ needs a lower
bound. \texttt{scipy.optimize.milp} returns \tablecmd{mip\_dual\_bound} and
\tablecmd{mip\_gap} (confirmed in SciPy~1.17)---a \emph{certified} lower bound on
each per-logical objective, valid even on timeout---but the current pipeline keeps
only the incumbent objective and the proven-optimal flag. In a future
non-early-exit pass, capturing $L_i=\lceil b_i\rceil$ for every relevant logical
subproblem, where $b_i$ is that subproblem's \tablecmd{mip\_dual\_bound}, gives
\[
  \min_i L_i \le d \le U,
\]
where $U$ is the best verified witness from MILP, \texttt{QDistEvol}, BP-OSD, or
another method. If every subproblem has a valid dual bound and $\min_i L_i=U$, the
distance is certified exact. If some subproblems are skipped or never started, the
global lower bound is still incomplete, so those cases should be reported as
partially bracketed rather than exact.

\paragraph{Closing more codes.} To actually prove exactness on borderline codes:
\begin{enumerate}
  \item \textbf{Use the heuristic witness as a cutoff target.} The current SciPy
  \texttt{milp} wrapper does not expose a MIP-start argument, but a verified
  \texttt{QDistEvol} witness still gives the target $U$: the MILP/SAT task becomes
  ``prove no logical of weight $<U$'' (or raise the dual bound to $U$). A backend
  with true MIP starts could additionally inject the witness as the initial
  incumbent.
  \item \textbf{Stronger solver:} if licensed, a direct Gurobi-backed run (e.g. a
  Webster/Gurobi method or our own Gurobi formulation) may be faster than
  HiGHS/SCIP on these instances.
  \item \textbf{SAT (\texttt{pySATDist}):} a different paradigm in which ``no
  logical of weight $<U$'' returning \textsc{unsat} is itself a certified lower
  bound; SAT sometimes closes instances MIP cannot.
  \item \textbf{CSS only:} a well-tuned BP-OSD (OSD-CS, more trials, DEM
  permutations) matched the recorded \texttt{ilp\_d} values in the tested
  overestimate set, often removing the need for MILP to find the CSS upper bound.
\end{enumerate}

\paragraph{Recommended workflow for the $n\ge108$ timeout codes.}
(i)~run a non-early-exit MILP pass that persists \tablecmd{mip\_dual\_bound} for
each attempted logical subproblem; (ii)~run \texttt{QDistEvol} to obtain a tight
independent $U$ (cheap; may beat the incumbent); (iii)~use $U$ as the cutoff target
for a closing MILP/SAT pass, optionally with Gurobi or \texttt{pySATDist};
(iv)~report $d\in[L,U]$, exact only where every relevant subproblem is bounded and
$L=U$, and flag the heuristic $U$ as ``probably tight'' otherwise.

\section{Limitations}

The cross-check is sound but bounded: it covers $n\le360$ and $d\le24$, and the
overestimate subset is low-distance. Heuristic tightness at larger $(n,k,d)$ is an
extrapolation until calibrated---on a held-out set of \emph{certified-exact} codes
one should measure the empirical probability that \texttt{QDistEvol} attains the
true distance as a function of $(n,k,d,\text{budget},\text{seed})$, and quantify
its run-to-run variance. Minimum distance is NP-hard (and hard to approximate), so
no method is both exact and scalable in general: beyond the regime where MIP or SAT
closes, the rigorous statement is the bracket $d\in[L,U]$ with a calibrated
confidence on $U$, not a point value. The structure of bicycle codes (circulant /
group-algebra), which black-box solvers discard, is the most promising route to
\emph{certified} statements at large scale.

\section{Conclusion}

An independent distance package corroborated our pipeline's recorded values on
198 codes with zero contradictions, including the up-to-$12\times$ BP-OSD
overestimate cases, and \texttt{QDistEvol} closed the BP-OSD gap against the
recorded MILP value on every one of the worst codes. For larger codes, the path
forward is not a magic exact oracle but a disciplined bracket: a cheap,
independently-verified heuristic upper bound, persisted MILP dual bounds, and
MILP/SAT closing passes for what remains closable.

\section*{Acknowledgments}

This report was drafted with the assistance of Anthropic's Claude Opus 4.8 and OpenAI GPT-5.5, which 
also helped implement and run the cross-check experiments described herein. All
results were verified against the committed data, and the author takes
responsibility for the content.

\begin{thebibliography}{9}
\bibitem{qcode} J.~Cruz-Benito, A.~W.~Cross, D.~Kremer, and I.~Faro,
\emph{Evolutionary discovery of bivariate bicycle codes with LLM-guided search},
arXiv:2606.02418 (2026).
\bibitem{webster} M.~Webster, A.~Jacob, and O.~Higgott,
\emph{Distance algorithms for classical codes, quantum codes and circuits},
arXiv:2603.22532 (2026).
\bibitem{lar} A.~J.~Landahl, J.~T.~Anderson, and P.~R.~Rice,
\emph{Fault-tolerant quantum computing with color codes}, arXiv:1108.5738 (2011).
\bibitem{bravyi} S.~Bravyi \emph{et al.},
\emph{High-threshold and low-overhead fault-tolerant quantum memory},
arXiv:2308.07915 (2023).
\end{thebibliography}

\end{document}
